\documentclass[11pt]{article}
\usepackage[margin=1in]{geometry}
\usepackage[T1]{fontenc}
\usepackage[utf8]{inputenc}
\usepackage{amsmath,amssymb,amsthm}
\usepackage{enumitem}

\title{ICPC World Finals 2023\\H. Jet Lag}
\author{}
\date{}

\begin{document}
\maketitle

\section*{Problem Summary}

We must choose sleep intervals so that no activity is missed. Sleeping for $k$ minutes gives $k$ more
minutes during which we are \emph{rested} and cannot fall asleep again, followed by another $k$ minutes
during which we may stay awake or fall asleep. We must start sleeping immediately at time $0$.

\section*{Key Observations}

\begin{itemize}[leftmargin=*]
    \item Ignoring the ``cannot sleep while rested'' restriction for a moment, the only useful sleep intervals
    are the free gaps
    \[
    [e_{i-1}, b_i], \qquad e_0 = 0.
    \]
    Sleeping any shorter interval inside the same gap can only decrease how far we can stay functional.
    \item If we sleep through a whole gap $[l,r]$, then the sleep lasts $r-l$ minutes, so afterwards we can
    stay functional until
    \[
    r + 2(r-l) = 3r - 2l.
    \]
    \item Thus, in the relaxed model, each gap acts like a transition $l \mapsto 3r-2l$. Scanning the gaps
    in time order and keeping every gap that increases the current reachable time computes the farthest
    reachable time.
    \item After this greedy selection, consecutive chosen sleeps can be shortened slightly so that the next
    sleep starts after the previous rested phase ends, while still keeping the next chosen start reachable.
\end{itemize}

\section*{Algorithm}

\begin{enumerate}[leftmargin=*]
    \item If the first activity starts at time $0$, output \texttt{impossible}, because we are required to sleep
    immediately.
    \item Let the free gaps be $[l_i,r_i]=[e_{i-1},b_i]$.
    \item Scan them from left to right, maintaining the farthest reachable time \texttt{reach} in the relaxed
    model:
    \begin{itemize}[leftmargin=*]
        \item if $l_i > \texttt{reach}$, the gap is unreachable;
        \item otherwise sleeping through it would extend reach to $3r_i-2l_i$;
        \item if this extension is larger than the current \texttt{reach}, keep the gap and update
        \texttt{reach}.
    \end{itemize}
    \item If the final \texttt{reach} is smaller than the end of the last activity, output \texttt{impossible}.
    \item Otherwise let the chosen relaxed sleeps be $[s_i,t_i]$. For each consecutive pair, shorten the
    earlier interval to
    \[
    t_i \leftarrow \min\!\left(t_i,\ \left\lfloor\frac{s_i+s_{i+1}}{2}\right\rfloor\right).
    \]
    Keep the last chosen interval unchanged.
    \item Output the shortened intervals.
\end{enumerate}

\section*{Correctness Proof}

We prove that the algorithm returns the correct answer.

\paragraph{Lemma 1.}
After processing the first $i$ gaps, the variable \texttt{reach} equals the farthest time reachable in the
relaxed model using only those gaps.

\paragraph{Proof.}
We use induction on $i$. Initially \texttt{reach}$=0$, which is optimal before using any gap.

Consider gap $[l_i,r_i]$. If $l_i > \texttt{reach}$, then this gap cannot be used by any schedule based on
the first $i$ gaps, so the optimum does not change. Otherwise the gap is reachable, and using it
completely extends the reachable time to $3r_i-2l_i$. Any relaxed schedule using only the first $i$ gaps
either skips this gap or uses it that way, so the new optimum is exactly
\[
\max(\texttt{reach},\, 3r_i-2l_i),
\]
which is the algorithm's update. \qed

\paragraph{Lemma 2.}
If the relaxed greedy phase reaches the end of the last activity, then the shortening step produces a valid
sleep schedule for the original problem.

\paragraph{Proof.}
Consider two consecutive chosen starts $s_i < s_{i+1}$. After shortening, we have
\[
t_i \le \frac{s_i+s_{i+1}}{2},
\]
so
\[
2t_i-s_i \le s_{i+1}.
\]
Thus the next sleep begins no earlier than the end of the previous rested phase.

At the same time, the original relaxed choice guaranteed that the next chosen start was reachable. Since
\[
\frac{s_i+s_{i+1}}{2} \ge \frac{2s_i+s_{i+1}}{3},
\]
the shortened endpoint still satisfies
\[
3t_i-2s_i \ge s_{i+1},
\]
so we can remain awake until the next chosen sleep begins. Therefore every shortened interval is valid,
and the last one still reaches far enough. \qed

\paragraph{Theorem.}
The algorithm outputs a valid schedule whenever one exists, and outputs \texttt{impossible} otherwise.

\paragraph{Proof.}
By Lemma 1, the greedy scan computes the maximum reachable time in the relaxed model. If even that
does not reach the end of the last activity, then no valid schedule exists in the original problem either.
If it does reach far enough, Lemma 2 converts the relaxed schedule into a valid original schedule.
Therefore the algorithm is correct. \qed

\section*{Complexity Analysis}

We scan the activities once and store the chosen intervals. The running time is $O(n)$ and the memory
usage is $O(n)$.

\section*{Implementation Notes}

\begin{itemize}[leftmargin=*]
    \item All times fit in 64-bit signed integers.
    \item The implementation verifies after shortening that each interval still has positive length and still
    reaches the next chosen start.
\end{itemize}

\end{document}
